% !TeX program = pdflatex

% ─────────────────────────────────────────────────────────────────────────────
% papers/mosfet.tex — Modelagem Dual Algébrica do MOSFET
%
% Objeto:
%   modelagem algébrica de um MOSFET por duas faces de realização:
%   formação do canal e transporte.
%
% Tese:
%   o campo de gate não é tratado como transporte de carga pelo gate;
%   ele realiza uma modulação da face de formação do canal.
%
%   O estrangulamento (pinch-off) é modelado como uma interface em que
%   a face de formação se anula localmente, sem que a face de transporte
%   seja anulada — exactamente como, no eletromagnetismo araniano,
%   estado local ≠ quantidade transportada.
%
% Arquitetura:
%   formação → transporte → interface → estrangulamento → saturação.
%
% Alinhamento com fisica.tex (Parte Eletromagnetismo):
%   o par (a,b) do MOSFET é lido como par de escritas local,
%   com forma, metal e cruzado no sentido da Def. fis:def:em-par
%   e da Obs. fis:obs:cisao. O duomorfismo é a mesma troca de face
%   já escrita no Corpo Universal.
%
% Estatuto:
%   modelagem algébrica. A confrontação quantitativa com dispositivos
%   físicos reais é etapa posterior (verificação da realização).
%
% Não é:
%   substituição da teoria de semicondutores;
%   identificação prematura entre objetos algébricos e grandezas físicas;
%   afirmação de equivalência física antes da verificação experimental.
%
% Compila:
%   pdflatex mosfet.tex
% ─────────────────────────────────────────────────────────────────────────────

\documentclass[11pt,a4paper]{article}

\usepackage[utf8]{inputenc}
\usepackage[T1]{fontenc}
\usepackage[brazil]{babel}

\usepackage{amsmath,amssymb,amsthm}

\usepackage[margin=2.4cm]{geometry}

\usepackage{booktabs}
\usepackage{array}
\usepackage{xcolor}

\usepackage[htt]{hyphenat}

\usepackage[hidelinks,breaklinks]{hyperref}

\sloppy
\emergencystretch=2em

\definecolor{cinza}{gray}{0.35}

\newcommand{\code}[1]{\texttt{\small #1}}

\newcommand{\caixa}[1]{%
  \par\smallskip\noindent
  \begin{center}
  \fbox{\begin{minipage}{0.94\textwidth}\centering\smallskip #1\smallskip\end{minipage}}
  \end{center}\smallskip}

\newcommand{\abs}[1]{\lvert #1\rvert}

\newcommand{\medido}[1]{%
  \par\medskip\noindent
  {\small\color{cinza}\textbf{Medido.} #1}
  \par\medskip}

\newcommand{\fis}[2]{\textbf{#2}~\textnormal{(\code{#1})}}

\providecommand{\pt}[1]{\textbf{P#1}}

\newcommand{\sepcol}{\quad\penalty0$\big|$\quad\penalty0}

\newcommand{\colunas}[3]{%
\par\smallskip\noindent{\footnotesize\raggedright
\textbf{prova}: #1\sepcol
\textbf{realiza}: #2\sepcol
\textbf{verifica}: #3\par}\smallskip}

\providecommand{\interfacen}[1]{}

\newtheorem{teorema}{Teorema}
\newtheorem{lema}[teorema]{Lema}
\newtheorem{definicao}[teorema]{Definição}
\newtheorem{corolario}[teorema]{Corolário}
\newtheorem{proposicao}[teorema]{Proposição}
\newtheorem{observacao}[teorema]{Observação}

\renewcommand{\proofname}{Demonstração}

% ─────────────────────────────────────────────────────────────────────────────
%  papers/gkcapa.tex — a capa do Reino, mínima, para os papers em `article`.
%
%  O estilo.tex completo é do `report` (títulos de capítulo, skak, …). Aqui fica
%  só o que a capa precisa, para o paper continuar a compilar sozinho com
%  pdflatex. O tradutor próprio (tests/tex) carrega o estilo.tex da raiz / o
%  symlink papers/estilo.tex e avalia o mesmo `\gkcapa`.
% ─────────────────────────────────────────────────────────────────────────────
\definecolor{ouro}{HTML}{B8912F}
\definecolor{ouroclaro}{HTML}{F3E9CE}
\definecolor{tinta}{HTML}{1E1B16}
\definecolor{regua}{HTML}{6E6858}
\providecommand{\gknota}{\fontsize{7.62}{11.05}\selectfont}
\providecommand{\gktexto}{\fontsize{10.50}{15.22}\selectfont}
\providecommand{\gksub}{\fontsize{12.33}{17.87}\selectfont}
\providecommand{\gksec}{\fontsize{14.47}{20.98}\selectfont}
\providecommand{\gkcap}{\fontsize{16.99}{24.63}\selectfont}
\providecommand{\gktit}{\fontsize{23.42}{33.95}\selectfont}
\providecommand{\Prj}{\mathbb{P}}
\providecommand{\Trf}{\mathcal{T}}
\providecommand{\gkcapa}[3]{%
\title{\vspace{-0.6cm}
{\color{ouro}\rule{\textwidth}{1.4pt}}\\[6mm]
\textbf{\gktit\color{tinta} Reino Dourado}\\[3mm]{\gksec\itshape\color{regua} Golden Kingdom}\\[8mm]
{\gkcap\scshape\color{ouro} #1}\\[5mm]
{\gksub\color{regua} #2}\\[2mm]
{\gktexto\color{regua}\itshape #3}\\[7mm]
{\color{ouro}\rule{\textwidth}{0.6pt}}\\[9mm]{\gknota\color{regua}\begin{minipage}{\textwidth}\setlength{\parindent}{0pt}%
\textbf{Aviso de Isenção Algorítmica:} O universo, as leis e as entidades contidas nestas páginas ---
incluindo felinos matriciais, aves de rapina algébricas e amuletos de controle de fluxo --- são
construções de pura ficção formal. Esta obra não descreve a realidade física, não serve como
engenharia reversa do cosmos e não constitui doutrina religiosa. Qualquer semelhança com bugs reais do seu
sistema operacional é mera coincidência (ou falta de reza). Prossiga por sua própria conta e
risco.\end{minipage}}}
\author{}\date{}}


\begin{document}

\interfacen{6}

\gkcapa
  {MOSFET --- modelagem dual algébrica}
  {formação do canal, transporte e estrangulamento}
  {par de escritas local alinhado ao eletromagnetismo araniano}

\maketitle

\begin{abstract}
\noindent
Este trabalho constrói uma modelagem algébrica dual para o MOSFET,
tomando como objeto a relação entre a formação do canal, o transporte
de carga e o estrangulamento (\emph{pinch-off}).

A construção parte da separação entre duas faces de realização:
uma face associada à disponibilidade do canal (determinada pela ação
do campo elétrico do gate) e uma face associada ao transporte ao longo
do canal. O ponto central é que o estrangulamento não constitui,
algebricamente, o desaparecimento do transporte: a face de formação
pode atingir valor nulo localmente enquanto a face de transporte
permanece realizada.

Esta leitura é alinhada à Parte de Eletromagnetismo de \code{fisica.tex}:
o par local do MOSFET é tratado como \textbf{par de escritas}, com
forma, metal e leitura orientada no sentido das Definições
\fis{fis:def:em-par}{par}, \fis{fis:def:em-forma}{forma} e
\fis{fis:def:em-fluxo}{cruzado}. A cláusula estrutural
\[
\boxed{\;\text{estado local}\;\neq\;\text{quantidade transportada}\;}
\]
é a mesma que separa, no eletromagnetismo araniano, a escrita local
da quantidade que atravessa a interface.

A hipótese é organizada pela estrutura dual do Corpo Universal:
duas faces, cisão, interface e duomorfismo. O objectivo desta etapa
é obter uma representação formal mínima da transição
\[
\text{canal} \longrightarrow
\text{estrangulamento} \longrightarrow
\text{transporte por campo},
\]
deixando a confrontação quantitativa com dispositivos físicos para
a etapa de verificação da realização.
\end{abstract}


\section{Objeto e estatuto da construção}

O objeto deste trabalho é o MOSFET enquanto sistema de controle
eléctrico de uma região de condução.

A construção não começa pela importação das equações usuais do
dispositivo. Começa pela pergunta estrutural:

\[
\boxed{
\text{como uma acção de campo pode modificar a disponibilidade
do transporte sem constituir, ela própria, o transporte?}
}
\]

No NMOS, a tensão \(V_{GS}\) modifica a concentração de portadores
na região superficial do semicondutor, enquanto \(V_{DS}\) estabelece
a queda de potencial ao longo do canal e, portanto, o campo responsável
pelo transporte.

A separação dessas duas funções fornece o primeiro par da modelagem:

\[
\boxed{
\text{formação}
\quad\big|\quad
\text{transporte}
}
\]

Esta separação é estrutural: não se identifica a grandeza que abre
o canal com a grandeza que conduz a carga. É a mesma recusa de fusão
que, no eletromagnetismo araniano, impede colapsar o par \((E,B)\)
numa única escrita ou promover o cruzado a operador \(\partial\).


\section{O par de escritas local do MOSFET}
\label{sec:par}

\begin{definicao}[par de escritas do MOSFET]
\label{def:par-mosfet}
Um \textbf{par de escritas local} do dispositivo é um par
\[
\mathcal M(x)=(a(x),b(x))
\]
sobre a mesma grade unidimensional do canal. \emph{Não é o campo \(G\)
de contagens da Mecânica}: \(G\) conta visitas; \(\mathcal M\) é uma
realização com duas componentes. Duas escritas \((\lambda a,\lambda b)\)
com \(\lambda\neq0\) pertencem à mesma classe na leitura local do metal
(Def.~\ref{def:metal-mosfet}).
\end{definicao}

A primeira componente representa a disponibilidade local do canal:

\[
a(x)=V_{GS}-V_{TH}-V_C(x),
\]

onde \(V_C(x)\) é o potencial local do canal.

A segunda componente representa a acção de transporte. Na primeira
realização tomamos

\[
b(x)=E_x(x),\qquad
E_x(x)=-\frac{dV_C}{dx}.
\]

Assim,

\[
\boxed{
\mathcal M(x)=
\left(
V_{GS}-V_{TH}-V_C(x),
-\frac{dV_C}{dx}
\right).
}
\]

A primeira face determina se a inversão pode permanecer realizada;
a segunda determina a acção de campo ao longo da região de transporte.
O par não deve ser colapsado numa única variável --- exactamente como
o par \((E,B)\) da Def.~\fis{fis:def:em-par}{par eletromagnético}
não se funde num único campo.


\section{Forma, metal e leitura orientada}
\label{sec:forma-metal}

Alinhando com a Parte de Eletromagnetismo, introduzimos as três leituras
do par local:

\begin{definicao}[forma local]
\label{def:forma-mosfet}
A \textbf{forma local} do par \(\mathcal M=(a,b)\) é
\[
\boxed{\;a^{2}-b^{2}\;}
\]
no sentido da Def.~\fis{fis:def:em-forma}{forma catalogada}. Não é um
corpo no sentido da Def.~\fis{fis:def:corpos}{corpos}: não é um troço
de \(I\), não se superpõe célula a célula. É leitura algébrica do
resíduo entre as duas faces.
\end{definicao}

\begin{definicao}[metal do par]
\label{def:metal-mosfet}
A \textbf{leitura local} do par, quando \(b\neq0\), é a razão
\[
\boxed{\;\frac{\lvert a\rvert}{\lvert b\rvert}\;}
\]
--- o \textbf{metal}, no sentido da Def.~\fis{fis:def:em-metal}{metal}.
O metal é o quociente; a interface de estrangulamento é o sítio onde
a face de formação se anula. São dois objectos.
\end{definicao}

\begin{definicao}[leitura orientada]
\label{def:cruzado-mosfet}
O cruzado do par, no canal unidimensional, é o produto orientado
\[
\boxed{\;a\cdot b\;}
\]
lido com o sinal da orientação do transporte. Regista a ordem na
cisão directo\,/\,cruzado já escrita (Obs.~\fis{fis:obs:cisao}{cisão}).
Não é o operador \(\partial\). \emph{Regista a orientação da passagem;
não é grandeza de estado.}
\end{definicao}

\begin{lema}[forma nula e metal]
\label{lem:forma-nula}
Se \(b\neq0\),
\[
\boxed{\;a^{2}-b^{2}=0\qquad\Longleftrightarrow\qquad\frac{\lvert a\rvert}{\lvert b\rvert}=1\;}
\]
--- a forma anula exactamente onde o metal vale um. Não é o vácuo do
canal; é a condição de anulação lida uma vez na forma e outra vez no
quociente (cf.\ Lema~\fis{fis:lem:em-cone}{forma nula e metal}).
\end{lema}

\begin{proof}
\(a^{2}-b^{2}=b^{2}\bigl((\lvert a\rvert/\lvert b\rvert)^{2}-1\bigr)\) e
\(b^{2}\neq0\).
\end{proof}


\section{Disponibilidade do canal}

A condição local de existência da camada de inversão é escrita,
nesta modelagem mínima, como

\[
a(x)>0.
\]

No limite,

\[
a(x)=0,
\]

a face de formação atinge a interface.

Para o extremo do drain,

\[
a(L)=V_{GS}-V_{TH}-V_{DS}.
\]

Consequentemente,

\[
a(L)=0
\]

equivale à condição

\[
\boxed{
V_{DS}=V_{GS}-V_{TH}.
}
\]

Esta igualdade será tomada como a condição algébrica de entrada
no estrangulamento.


\section{Estrangulamento como interface}
\label{sec:interface}

O ponto central da construção é não identificar

\[
a(L)=0
\]

com

\[
\mathcal M(L)=(0,0).
\]

No estrangulamento,

\[
\boxed{
\mathcal M(L)=(0,b(L)),
}
\]

com

\[
b(L)\neq0.
\]

Portanto, o que desaparece é a face de formação do canal naquele
ponto; a face de transporte permanece.

A leitura estrutural é:

\[
\boxed{
\text{formação}\;\longrightarrow\;
\text{interface}\;\longrightarrow\;
\text{transporte}.
}
\]

O estrangulamento é, portanto, uma mudança de realização local,
e não a anulação do objecto transportado. É a realização concreta,
no dispositivo, da cláusula

\[
\boxed{\;\text{estado local}\;\neq\;\text{quantidade transportada}\;}
\]

já enunciada na tese do Eletromagnetismo araniano
(\S\fis{fis:em-tese}{tese EM}).


\section{Campo e persistência do transporte}

A persistência do transporte após o estrangulamento é representada
pela presença de campo eléctrico na região de interface:

\[
\mathbf F=q\mathbf E.
\]

A ausência local de uma camada de inversão não implica ausência
de força sobre os portadores.

Na representação dual,

\[
a(L)=0
\qquad\text{enquanto}\qquad
b(L)\neq0.
\]

Assim, a passagem não é representada por uma continuidade uniforme
do canal. Ela é representada por uma travessia entre duas faces:

\[
\boxed{
(a,b)
\;\longrightarrow\;
(0,b')
}
\]

com \(b'\) determinado pela realização de transporte.


\section{Duomorfismo e troca de leitura}
\label{sec:duomorf}

A estrutura anterior permite introduzir o duomorfismo como a
interface entre as duas faces, no sentido do Teor.~\fis{fis:thm:em-hodge}{Hodge}
e da Obs.~\fis{fis:obs:cisao}{cisão}.

Se

\[
D:\oplus\longleftrightarrow\otimes
\]

é a troca de realização entre as operações directa e cruzada,
a leitura do MOSFET utiliza a mesma estrutura sem identificar
as duas faces.

Esquematicamente:

\[
\boxed{
\text{formação}
\overset{D}{\longleftrightarrow}
\text{transporte}.
}
\]

A função do duomorfismo não é afirmar que formação e transporte
são a mesma operação. Ao contrário: ele regista a passagem
reversível entre as duas leituras do mesmo suporte.

O estrangulamento é então localizado na interface dessa passagem.


\section{A dobra no estrangulamento}

Se \(\partial\) representa a dobra da realização, o pinch-off
pode ser lido como uma mudança de representação:

\[
\partial:
(a,b)
\longrightarrow
(0,b').
\]

A propriedade essencial não é a anulação completa do estado,
mas a preservação da travessia sob mudança de face.

Em particular,

\[
\boxed{
\text{estrangulamento}
\neq
\text{bloqueio}.
}
\]

O estrangulamento marca o ponto em que a primeira representação
do transporte deixa de ser válida localmente e a realização passa
a ser descrita pela região de campo.


\section{Regime de saturação}

A condição

\[
V_{DS}\geq V_{GS}-V_{TH}
\]

define, na teoria convencional de canal longo, a entrada na região
de saturação.

Na presente leitura, essa condição será interpretada como a
condição em que a face de formação não pode mais ser estendida
uniformemente até o drain.

A queda adicional de potencial passa a ser realizada na região
de estrangulamento.

Assim, a saturação não é introduzida como uma regra independente.
Ela deverá emergir da estrutura:

\[
\boxed{
\text{formação}
\rightarrow
\text{deformação}
\rightarrow
\text{interface}
\rightarrow
\text{saturação}.
}
\]


\section{Construção da corrente a partir do par}
\label{sec:corrente}

A etapa seguinte consiste em construir a corrente a partir do par

\[
\mathcal M(x)=(a(x),b(x))
\]

sem importar a fórmula de saturação como axioma.

\begin{proposicao}[densidade de corrente local]
\label{prop:densidade}
Na região onde \(a(x)>0\), a densidade de corrente longitudinal
é proporcional ao produto das faces:
\[
\boxed{\;j(x)\;\propto\;a(x)\,b(x)\;}
\]
--- o cruzado do par (Def.~\ref{def:cruzado-mosfet}). Na interface
\(a=0\) a contribuição da face de formação anula-se localmente,
mas a face de transporte \(b\) permanece e a corrente continua
a ser realizada pela região de campo.
\end{proposicao}

\begin{observacao}
A proporcionalidade \(j\propto a\cdot b\) é a realização mínima
da leitura orientada. Não se identifica ainda \(a\) com carga de
inversão nem \(b\) com mobilidade\(\times\)campo; essas
identificações pertencem à etapa de verificação experimental.
\end{observacao}

Integrando ao longo do canal, no regime em que a face de formação
se mantém positiva até ao drain (\(V_{DS}<V_{GS}-V_{TH}\)),
obtém-se a dependência linear-quadrática clássica. No regime de
saturação a integração detém-se na interface \(a(L_{\mathrm{eff}})=0\)
e a corrente deixa de crescer com \(V_{DS}\).

O alvo de verificação é recuperar

\[
I_D
=
\frac12 k_n
\bigl(V_{GS}-V_{TH}\bigr)^2
\]

no regime de saturação, como \emph{consequência} da integração do
par dual, e não como ponto de partida.

O critério permanece:

\[
\boxed{
\text{álgebra dual}
\;\Longrightarrow\;
\text{lei de transporte}
\;\Longrightarrow\;
\text{equação MOSFET}
}
\]

e não o caminho inverso.


\section{Mapa de alinhamento com o Eletromagnetismo}
\label{sec:mapa}

\begin{center}
\small
\begin{tabular}{@{}p{0.32\textwidth}p{0.32\textwidth}p{0.28\textwidth}@{}}
\toprule
\textbf{MOSFET} & \textbf{Eletromagnetismo (\code{fisica.tex})} & \textbf{estatuto} \\
\midrule
par \(\mathcal M=(a,b)\) & par de escritas \((E,B)\) & Def.~\ref{def:par-mosfet} / \fis{fis:def:em-par}{par} \\
forma \(a^{2}-b^{2}\) & forma \(E^{2}-B^{2}\) & Def.~\ref{def:forma-mosfet} / \fis{fis:def:em-forma}{forma} \\
metal \(\lvert a\rvert/\lvert b\rvert\) & metal \(\lvert E\rvert/\lvert B\rvert\) & Def.~\ref{def:metal-mosfet} / \fis{fis:def:em-metal}{metal} \\
cruzado \(a\cdot b\) & cruzado \(E\times B\) & Def.~\ref{def:cruzado-mosfet} / \fis{fis:def:em-fluxo}{fluxo} \\
interface \(a=0\), \(b\neq0\) & estado local \(\neq\) quantidade transportada & \S\ref{sec:interface} / \S\fis{fis:em-tese}{tese} \\
duomorfismo \(D\) & Hodge / \(\tau\) no plano & \S\ref{sec:duomorf} / Teor.~\fis{fis:thm:em-hodge}{Hodge} \\
\bottomrule
\end{tabular}
\end{center}

\medskip\noindent
\emph{Nenhuma identificação física prematura.} O mapa é estrutural:
as mesmas dualidades, a mesma recusa de fusão, a mesma preservação
da travessia sob mudança de face.


\section{Próximas realizações}

\begin{enumerate}
\item Explicitar a integração do produto \(a(x)b(x)\) sob as duas
  condições de contorno (linear e saturação) e obter as expressões
  de \(I_D\) como teoremas de realização.
\item Introduzir a mobilidade e a capacitância de óxido apenas como
  factores de escala da realização, nunca como axiomas da dualidade.
\item Confrontar numericamente com o modelo de nível~1 de SPICE
  (canal longo) e registar os desvios como medidores de verificação.
\item Estender o par \(\mathcal M\) a duas dimensões (largura do canal)
  mantendo a separação formação\,/\,transporte.
\end{enumerate}

\end{document}
